\chapter{Design Specification and Current System Architecture}
\label{ch:design}

%==============================================================================
\section{Data Pipeline and Representation}
%==============================================================================

\subsection{Source and Canonical Input}

The operational pipeline starts from hourly ICU tables derived from MIMIC-III \cite{johnson2016mimic}
through a MIMIC-Extract style workflow \cite{wang2020mimicextract}. The canonical training input for autoregressive modelling is \code{data/processed/autoregressive\_data.csv}.
Hour-0 training uses \code{data/processed/hour0\_data.csv}.

\subsection{Autoregressive Transformation}

Let $\mathcal{S}$ be static columns and $\mathcal{X}$ be dynamic columns.
With lag-1 engineering, each row at time $t$ uses:
\begin{align}
\text{target}_t &= \mathbf{x}_t \in \mathbb{R}^{|\mathcal{X}|}, \\
\text{condition}_t &= [\mathbf{s}, \mathbf{x}_{t-1}] \in \mathbb{R}^{|\mathcal{S}|+|\mathcal{X}|}.
\end{align}

\paragraph{Lag-0 and missing lag cells (implemented).}
The autoregressive table is built by sorting on subject identifier and hour index and, for each dynamic column, producing a lag-1 feature by shifting within each subject group.
\textbf{Lag-0 (no in-group history).} The first row of each subject has no prior row in that group, so the lag feature is undefined before filling.
\textbf{Fill strategy and value.} Undefined lag entries are replaced by the float sentinel $-1.0$, yielding a dense, rectangular lag feature matrix (see Implementation chapter for the concrete pipeline step and parameters).
The same fill applies to missing entries arising when $t-1$ exists but the source value at $t-1$ was missing: the shift propagates the missing marker, which is then indistinguishable from lag-0 after filling.

\paragraph{Rationale and risks.}
The sentinel keeps a dense, rectangular lag feature matrix so hour~0 rows remain in the autoregressive CSV and need not be dropped for lack of history; the IID hour-0 pipeline can still model the initial state separately.
Risks: \code{-1.0} is not a reserved missing code in the raw feature space and it conflates ``no prior timestep'' with ``prior value missing.'' Chapter~7 links this design simplification to observed fidelity limitations.

\subsection{Column Roles Persisted as Metadata}

The preprocessor artifact stores:
\begin{itemize}\tightlist
    \item \code{target\_cols} (semantic generation outputs),
    \item \code{condition\_cols} (semantic conditioning inputs),
    \item \code{target\_indices} and \code{condition\_indices} in transformed space.
\end{itemize}

This design persists both semantic columns and transformed indices because transformed column order is type-grouped (numeric, binary, categorical), not semantic target-then-condition order.

%==============================================================================
\section{Tabular Preprocessor Design}
%==============================================================================

\subsection{Feature-Type Identification}

Feature typing is centralised in a dedicated feature-typing utility and consumed by the preprocessor fit step.
Columns are grouped into:
\begin{itemize}\tightlist
    \item numeric continuous,
    \item binary 0/1,
    \item categorical,
    \item integer-like numeric (for inverse rounding policy).
\end{itemize}

\subsection{Transform Semantics (Current)}

The current \code{TabularPreprocessor} is intentionally simple:
\begin{enumerate}\tightlist
    \item Numeric: pass-through (no normalization).
    \item Binary: pass-through in transform; round/clip policies in inverse path.
    \item Categorical: label encoding to integer IDs.
\end{enumerate}

This design differs from one-hot plus min-max encoding variants common in tabular generative baselines and aligns with the tree-based backend, where strict scaling is not required for split-based learners \cite{chen2016xgboost}.

\subsection{Streaming Fit for Full-Dataset Coverage}

The preprocessor supports streaming chunked fitting over the full CSV, avoiding a materialised in-memory copy of the training table.

The design goal is to capture rare categorical values and robust numeric ranges before any subsampled training.

\subsection{Inverse Transform Contract}

\code{inverse\_transform()} reconstructs a human-readable DataFrame by:
\begin{itemize}\tightlist
    \item decoding categorical IDs back to labels,
    \item applying binary/integer rounding rules,
    \item preserving numeric values in raw scale.
\end{itemize}

This decode step is part of the intended export path for downstream CSV outputs and clinical interpretability, even though model training happens in encoded numeric space.

\subsection{Index-Based Split as a Design Constraint}

Because transform output order is type-grouped:
\begin{equation}
\text{transformed\_order} = [\text{numeric},\text{binary},\text{categorical}],
\end{equation}
target and condition blocks may become non-contiguous in transformed coordinates.

Therefore training scripts use:
\begin{equation}
X_{\text{target}} = X[:, \text{target\_indices}], \quad
X_{\text{condition}} = X[:, \text{condition\_indices}],
\end{equation}
instead of naive prefix slicing.

%==============================================================================
\section{Generator Backbone Architecture}
%==============================================================================

\subsection{System Pipeline Overview}

Figure~\ref{fig:pipeline-overview} shows the end-to-end pipeline from raw MIMIC-III ICU tables to evaluated synthetic output. The two-stage architecture is explicit: the hour-0 IID model and the autoregressive transition model are trained independently on their respective data slices, and synthetic trajectories are assembled by chaining their outputs.

% , , , , , , , , , , , , , , , , , , , , , , , , ,
% Figure 1: Two-stage pipeline overview
% , , , , , , , , , , , , , , , , , , , , , , , , ,
\begin{figure}[htbp]
\centering
\begin{tikzpicture}[scale=0.88, transform shape,
  node distance = 6mm and 12mm,
  box/.style = {rectangle, draw, rounded corners=2pt,
                minimum width=32mm, minimum height=8mm,
                text centered, font=\small},
  data/.style = {box, fill=blue!10},
  proc/.style = {box, fill=orange!15},
  model/.style= {box, fill=green!15},
  eval/.style = {box, fill=red!10},
  arr/.style  = {-Stealth, thick},
  grp/.style  = {draw=gray!60, dashed, rounded corners=4pt,
                 inner sep=4pt, fill=gray!5}
]

% ,  Data preparation column (left) ,
\node[data]  (mimic)   {MIMIC-III ICU tables};
\node[proc,  below=of mimic]  (p00)    {Prepare hour-0 table};
\node[proc,  below=of p00]    (p01)    {Prepare autoregressive table};
\node[proc,  below=of p01]    (p01b)   {Fit full-data preprocessor};

% ,  Hour-0 branch (middle column) ,
% Anchored at the height of p00 so both training nodes are at the same level
\node[model, right=24mm of p00]  (h0train)  {Train hour-0 IID model};
\node[model, below=of h0train]   (h0gen)    {Sample initial states $\hat{\mathbf{x}}_0$};

% ,  Autoregressive branch (right column, same top height as h0 branch) ,
\node[model, right=20mm of h0train] (artrain)  {Train AR model};
\node[model, below=of artrain]      (argen)    {Autoregressively roll-out $\hat{\mathbf{x}}_{1:T}$};

% ,  Merge: centred between h0gen and argen ,
\node[proc] at ($(h0gen.south)!0.5!(argen.south) + (0,-16mm)$) (merge)
      {Assemble synthetic cohort CSV};

% ,  Evaluation stack (below merge) ,
\node[eval, below=of merge] (tstr)   {Patient-level TSTR (mortality, LOS)};
\node[eval, below=of tstr]  (qual)   {KS / Correlation / Range violations};
\node[eval, below=of qual]  (priv)   {DCR + DOMIAS-style MIA};

% ,  Arrows: data prep ,
\draw[arr] (mimic) -- (p00);
\draw[arr] (p00)   -- (p01);
\draw[arr] (p01)   -- (p01b);

% ,  Arrows: data prep → branches ,
\draw[arr] (p00)   -- (h0train);
\draw[arr] (p01b)  -- (artrain);

% ,  Arrows: within branches ,
\draw[arr] (h0train) -- (h0gen);
\draw[arr] (artrain) -- (argen);

% ,  Arrows: converge to merge ,
\draw[arr] (h0gen.south) -- (merge.north west);
\draw[arr] (argen.south) -- (merge.north east);

% ,  Arrows: evaluation ,
\draw[arr] (merge) -- (tstr);
\draw[arr] (tstr)  -- (qual);
\draw[arr] (qual)  -- (priv);

% ,  Grouping boxes ,
\begin{scope}[on background layer]
  \node[grp, fit=(p00)(p01)(p01b), label=above:{\scriptsize Data preparation}] {};
  \node[grp, fit=(h0train)(h0gen), label=above:{\scriptsize Hour-0 model}] {};
  \node[grp, fit=(artrain)(argen), label=above:{\scriptsize Autoregressive model}] {};
  \node[grp, fit=(tstr)(qual)(priv), label=below:{\scriptsize Evaluation stack}] {};
\end{scope}

\end{tikzpicture}
\caption{End-to-end pipeline. MIMIC-III ICU tables feed two parallel training paths: an IID hour-0 model (generating initial patient states) and an autoregressive transition model (generating subsequent hourly observations). Their outputs are chained into a synthetic cohort CSV, which is passed through the three-axis evaluation stack.}
\label{fig:pipeline-overview}
\end{figure}

\subsection{ForestFlow Architecture}

Figure~\ref{fig:forestflow-arch} shows how ForestFlow is instantiated in this system. The formal flow-matching equations are defined in Chapter~4; here the design point is implementation-facing: one regressor bank per time level, clean conditioning, and Euler rollout for generation.

% , , , , , , , , , , , , , , , , , , , , , , , , ,
% Figure 2: ForestFlow backbone
% , , , , , , , , , , , , , , , , , , , , , , , , ,
\begin{figure}[htbp]
\centering
\begin{tikzpicture}[scale=0.88, transform shape,
  node distance = 5mm and 12mm,
  box/.style  = {rectangle, draw, rounded corners=2pt,
                 minimum width=30mm, minimum height=8mm,
                 text centered, font=\small},
  blue/.style = {box, fill=blue!10},
  orange/.style={box, fill=orange!15},
  green/.style ={box, fill=green!15},
  red/.style   ={box, fill=red!10},
  arr/.style   = {-Stealth, thick},
  darr/.style  = {-Stealth, thick, dashed},
  grp/.style   = {draw=gray!60, dashed, rounded corners=4pt,
                  inner sep=4pt, fill=gray!5}
]

% , - Training path (left column) , -
\node[blue]   (noise)  {Sample $\mathbf{z}\sim\mathcal{N}(0,I)$};
\node[blue,   below=of noise]   (interp) {Interpolate: $\mathbf{x}_t=(1{-}t)\mathbf{z}+t\mathbf{x}$};
\node[orange, below=of interp]  (concat) {Concat $[\mathbf{x}_t^{\text{noisy}} \mid \mathbf{c}_{\text{clean}}]$};
\node[green,  below=of concat]  (xgb)    {XGBoost regressor fits $\hat{\mathbf{v}}_{t_k}$};
\node[orange, below=of xgb]     (loss)   {Loss: $\|\hat{\mathbf{v}}_{t_k}-(\mathbf{x}-\mathbf{z})\|_2^2$};

% , - Time loop annotation , -
\node[right=5mm of xgb, font=\scriptsize, text=gray] (timeann)
      {$\times\; n_t$ time levels};

% , - Sampling path (right column) , -
\node[blue,   right=28mm of noise]   (znoise) {Start: $\hat{\mathbf{x}}_0\sim\mathcal{N}(0,I)$};
\node[green,  below=of znoise]       (pred)   {Predict $\hat{\mathbf{v}}_{t_k}(\hat{\mathbf{x}}_k,\mathbf{c})$};
\node[orange, below=of pred]         (euler)  {Euler step: $\hat{\mathbf{x}}_{k+1}=\hat{\mathbf{x}}_k+h\hat{\mathbf{v}}_{t_k}$};
\node[blue,   below=of euler]        (out)    {Output: $\hat{\mathbf{x}}_{n_t}$ (decoded)};

% , - Arrows: training , -
\draw[arr] (noise)  -- (interp);
\draw[arr] (interp) -- (concat);
\draw[arr] (concat) -- (xgb);
\draw[arr] (xgb)    -- (loss);

% , - Arrows: sampling , -
\draw[arr] (znoise) -- (pred);
\draw[arr] (pred)   -- (euler);
\draw[arr] (euler)  to[out=0,in=0,looseness=2]
           node[right,font=\scriptsize]{$k{=}0{\to}n_t{-}1$} (pred);
\draw[arr] (euler)  -- (out);

% , - Training → Sampling arrow , -
\draw[darr] (xgb.east) -- node[above,font=\scriptsize]{stored regressors} (pred.west);

% , - Group boxes , -
\begin{scope}[on background layer]
  \node[grp, fit=(noise)(interp)(concat)(xgb)(loss),
        label=above:{\scriptsize Training}] {};
  \node[grp, fit=(znoise)(pred)(euler)(out),
        label=above:{\scriptsize Sampling (Euler)}] {};
\end{scope}

\end{tikzpicture}
\caption{ForestFlow backbone. \textit{Training}: for each time level $t_k$, the target is interpolated with Gaussian noise and a regressor is fitted to predict the flow direction. \textit{Sampling}: starting from noise, $n_t$ first-order Euler steps integrate the learned vector field back to data space. All target dimensions are treated as continuous flows.}
\label{fig:forestflow-arch}
\end{figure}

\subsection{HS3F Architecture}

Figure~\ref{fig:hs3f-arch} shows the HS3F system design. Chapter~4 defines the formal heterogeneous factorisation; this section focuses on architecture: type-routed heads, feature-wise left-to-right generation, and explicit prefix feedback during sampling.

% , , , , , , , , , , , , , , , , , , , , , , , , ,
% Figure 3: HS3F backbone
% , , , , , , , , , , , , , , , , , , , , , , , , ,
\begin{figure}[htbp]
\centering
\begin{tikzpicture}[scale=0.78, transform shape,
  node distance = 5mm and 10mm,
  box/.style   = {rectangle, draw, rounded corners=2pt,
                  minimum width=30mm, minimum height=8mm,
                  align=center, font=\small},
  blue/.style  = {box, fill=blue!10},
  orange/.style= {box, fill=orange!15},
  green/.style = {box, fill=green!15},
  purple/.style= {box, fill=purple!10},
  red/.style   = {box, fill=red!10},
  arr/.style   = {-Stealth, thick},
  darr/.style  = {-Stealth, thick, dashed},
  grp/.style   = {draw=gray!60, dashed, rounded corners=4pt,
                  inner sep=4pt, fill=gray!5}
]

% , - Shared input , -
\node[blue]   (cond)  {Condition $\mathbf{c} = [\mathbf{s},\, \mathbf{x}_{t-1}]$};

% , - Feature loop head (router) , -
\node[orange, below=10mm of cond]  (router) {Feature $j$ router: type check};
\node[right=3mm of router, font=\scriptsize, text=gray]
     (loopann) {for $j=1,\ldots,d_x$};

% , - Continuous branch , -
\node[green, below left=12mm and 12mm of router]  (flow)
     {Flow-matching regressor\\$\hat{v}_{t_k}(x_{t,j}\mid\mathbf{c},x_{t,<j})$};
\node[green, below=of flow] (rk4)
     {RK4 integration $\to \hat{x}_{t,j}$};

% , - Categorical branch , -
\node[purple, below right=12mm and 12mm of router] (clf)
     {XGBoost classifier\\$\hat{p}(x_{t,j}\mid\mathbf{c},x_{t,<j})$};
\node[purple, below=of clf] (sample)
     {Sample $\hat{x}_{t,j}\sim\hat{p}(\cdot)$};

% , - Accumulate , -
\node[blue, below=20mm of router]  (accum)
     {Append $\hat{x}_{t,j}$ to generated prefix};
\node[orange, below=of accum] (done)
     {All $d_x$ features generated: output $\hat{\mathbf{x}}_t$};

% , - Arrows , -
\draw[arr] (cond)   -- (router);
\draw[arr] (router) -- node[left, font=\scriptsize]{continuous} (flow);
\draw[arr] (router) -- node[right,font=\scriptsize]{categorical} (clf);
\draw[arr] (flow)   -- (rk4);
\draw[arr] (clf)    -- (sample);
\draw[arr] (rk4)    -- (accum);
\draw[arr] (sample) -- (accum);
\draw[arr] (accum)  -- (done);
% Feedback: generated feature feeds back into condition for next j
\draw[darr] (accum.west) to[out=180,in=180,looseness=1.4]
            node[left,font=\scriptsize]{feeds $x_{t,<j+1}$} (router.west);

% , - Group boxes , -
\begin{scope}[on background layer]
  \node[grp, fit=(flow)(rk4), label=left:{\scriptsize Continuous head}] {};
  \node[grp, fit=(clf)(sample),  label=right:{\scriptsize Categorical head}] {};
\end{scope}

\end{tikzpicture}
\caption{HS3F backbone. Features are generated sequentially in a fixed order. Continuous features are produced by a flow-matching regressor integrated with RK4 by default (explicit Euler is also selectable at sampling time); categorical features are produced by an XGBoost classifier with multinomial sampling. Each generated feature is appended to the conditioning prefix before the next feature is generated, capturing inter-feature dependencies.}
\label{fig:hs3f-arch}
\end{figure}

\subsection{Backbone Selection}

Backbone selection is exposed at training time, with HS3F as the default; operational invocation details are documented in the Implementation chapter.

\subsection{ForestFlow Design Interface}

ForestFlow consumes a target matrix, a clean condition matrix, and feature-type metadata, and emits synthetic targets in transformed coordinates.
The implementation-level rule is simple: all target dimensions are routed through continuous flow regressors to keep interpolation arithmetic well defined.
This keeps the training path uniform but can reduce categorical fidelity; that risk is evaluated in Chapter~7.

\subsection{HS3F Design Interface}

HS3F consumes the same target/condition matrices plus feature typing, and emits synthetic targets in transformed coordinates.
Its implementation rule is to split continuous and categorical dimensions into separate heads and generate features left-to-right with prefix feedback.
Label remapping is retained so classifier class indices stay consistent under subsampling; this prevents index mismatch failures during sampling.

\subsection{Hour-0 + Autoregressive Coupling}

The architecture is intentionally two-stage:
\begin{enumerate}\tightlist
    \item Hour-0 IID model learns initial states.
    \item Autoregressive model learns transitions conditioned on lagged dynamics and static covariates.
\end{enumerate}

This decomposition separates structural initialization from transition dynamics more explicitly than a single model that must infer both jointly.

\paragraph{Sweep-level design.}
Sweep-level design, the hyperparameter protocol, transformed-matrix caching, and canonical result schema, is part of the operational pipeline rather than the generator architecture, and is specified in the Implementation chapter (see~\ref{sec:unified-sweep}).

%==============================================================================
\section{Design Rationale and Trade-offs}
%==============================================================================

\subsection{Why Raw Numeric Pass-Through?}

Tree boosting is split-based and generally robust to monotonic scaling transformations \cite{chen2016xgboost}.
Removing mandatory min-max scaling is intended to simplify preprocessing and reduce opportunities for inverse-transform coupling errors.

\subsection{Why Keep Inverse Transform?}

Even with simplified encoding, inverse transformation remains important for:
\begin{itemize}\tightlist
    \item readable synthetic table export,
    \item downstream evaluator compatibility,
    \item sanity checks in clinical units.
\end{itemize}

\subsection{Why Multi-Backbone Support?}

Forest-Flow \cite{jolicoeur2023forest} remains a well-established baseline and compatibility path.
HS3F \cite{akazan2024hs3f} is the default backbone; its literature positioning is reviewed in Chapter~\ref{ch:background}.

\subsection{Why Preserve Explicit Metadata Artifacts?}

Persisting split metadata and model type in pickled artifacts enables:
\begin{itemize}\tightlist
    \item backward-compatible loading in generation scripts,
    \item reproducible sweep reruns,
    \item traceable migration between architecture revisions.
\end{itemize}

Having justified the system architecture and design trade-offs, the next chapter documents how these design choices were implemented and made reproducible.
